\subsubsection{Hardware Revision Defects}
Three hardware defects were identified in this revision:
\begin{itemize}
	\item \textbf{Gate 03 input (FR-03):} The assigned pin PC9 does not support
	EXTI interrupts simultaneously with the other assigned interrupt lines.
	A bodge wire was soldered directly between PC9 and PC6 on the MCU pins,
	relocating the signal to a pin with full EXTI support, and the firmware was
	updated accordingly (Figure~\ref{fig:gpio-bodge}).
	A PCB revision with the correct pin assignment is required to eliminate the bodge.
	\item \textbf{Slide potentiometer LEDs (FR-04):} PWM-capable timer channels
	for the fader LEDs were incorrectly assigned during schematic entry, leaving
	the slider LED indicators non-functional in this revision.
	\item \textbf{BAT54S audio input clamp (EC-03):} A mismatch between the SPICE
	simulation model and the actual SOT-23 package pinout of the BAT54S dual Schottky
	diode caused the component to be placed with reversed orientation on the PCB.
	The error was corrected by bending the component pins upward and soldering it
	inverted, restoring the correct diode orientation (Figure~\ref{fig:diode-bodge}).
	A corrected footprint assignment is required in the next PCB revision.
\end{itemize}

\subsubsection{Slice Selection in Reverse Playback}
Slice boundary handling during reverse playback is not fully implemented; when reverse mode is active, slice selection is ignored and the full buffer is played back in reverse instead.

\subsubsection{Memory Headroom}
DTCMRAM utilisation stands at 98.5\% and main SRAM at 93\% in the current firmware build.
While sufficient for the implemented feature set, this leaves limited headroom for additional DSP algorithms or larger sample buffers.
The current recording and playback buffer holds \qty{2.5}{\second} of audio at $f_s = \qty{48}{\kilo\hertz}$, which is adequate for the tape-style use case but restrictive for longer recordings.

\subsubsection{Noise Characterisation (EC-05)}
Quantitative noise floor measurements have not been completed for this revision.
The analogue/digital separation strategy was applied during layout, but its effectiveness has not been formally verified through measurement.
During in-system testing in a Eurorack environment, the noise performance was subjectively sufficient and no audible noise floor was observed.

\subsubsection{Fixed-Point Phase Accumulator}
The playback engine uses a Q48.16 fixed-point phase accumulator, motivated by the concern that floating-point accumulation errors could cause phase drift over long playback durations (Section~\ref{sec:numerical-architecture}).
In retrospect, this concern does not apply to this system: the sample buffer holds only \qty{2.5}{\second} of audio at $f_s = \qty{48}{\kilo\hertz}$, meaning the playhead wraps within $2.5 \times 48000 = 120{,}000$ samples.
Floating-point rounding errors are on the order of $\epsilon_{\text{mach}} = 2^{-23} \approx 1.19 \times 10^{-7}$ per accumulation step for IEEE~754 single-precision arithmetic \cite[Table~3.5]{IEEEStandardFloatingPoint2019}; over the maximum 120\,000 steps before buffer wrap, the accumulated error remains below $0.012$ samples, well below the resolution of the interpolator.
A \texttt{float} playhead would be fully sufficient and would eliminate the fixed/floating boundary at the interpolator input, improving code readability. 
The fixed-point design is correct but unnecessary for this buffer size, and could be simplified in a future refactor.

\subsubsection{Input Stage Low-Frequency Rolloff Discrepancy}
The measured \SI{-3}{\decibel} cutoff frequency of the audio input high-pass filter is approximately \SI{37.6}{\Hz}, compared to a theoretical value of \SI{1.8}{\Hz} calculated from the AC-coupling network ($R_{in} = \SI{40}{\kilo\ohm}$, $C_3 = \SI{2.2}{\micro\farad}$).

The root cause has not been identified within the scope of this project, as only an end-to-end frequency response measurement was conducted.
To isolate the source, individual frequency response measurements of the input and output stages should be performed separately in a future revision.
Practically, the measured cutoff is reaching into the audible low frequency end, as sub-bass frequencies do reach the lower \qty{30}{Hz}.
Still, this rolloff is marginal and the device is usable for musical applications.